
\addchapnonumber{B\ \ Nouvelle: "Second souffle, Mirage ou virage ?"} \label{annex_nouvelle}


%\chapter{Nouvelle: \\ "Second souffle, \\ Mirage ou virage ?"}\label{annex_nouvelle}
%\chaptermark{\color{curcolor} Nouvelle: "Second souffle, Mirage ou virage ?"}




\poetry{6}{
Je cours pour être le meilleur...

\hspace*{2.5cm} ...de moi-même.

J'écris pour partager le meilleur...

\hspace*{2.5cm} ...de ma plume.

J'aime pour vivre le meilleur...

\hspace*{2.5cm} ...de ma vie.

}


\newpage



\includepdf[scale=1.05,pages=-]{others/Trail.pdf}



\newpage


\vspace*{\fill}

{\noindent\Large\textbf{Présentation}}
\\

Au cours de ma dernière année de thèse, j'ai participé à l'atelier d'écriture créative sur le climat organisé à l'ENS-PSL par la climatologue Aglaé Jézéquel (LMD-ENS) et l'écrivain Mathieu Simonet. Cet atelier convie chacun.e des participant.e.s à écrire une nouvelle avec un message autour du climat. Les nouvelles ainsi rédigées sont toutes disponibles sur le site de l'IPSL (\url{https://www.ipsl.fr/decouvrir/mozaika/chroniques-du-changement-climatique/}; dernier accès : Septembre 2023). Je tenais à ce que cette nouvelle fasse partie intégrante de mon manuscrit de thèse, car elle est à l'image de mon engagement sur la question climatique, à travers une approche sensible cette fois. Je tiens notamment à remercier les hydroclimatologues Florence Habets (LG-ENS) et Agnès Ducharne (METIS), ainsi que Paul Bonhomme (guide de haute montagne) pour les échanges que nous avons pu avoir au début de ce projet sur le cycle de l'eau en montagne face au changement climatique, mais aussi sur l'effort d'endurance.

\begin{spacing}{5}
~
\end{spacing}


{\noindent\Large\textbf{Éléments de contexte}}
\\

Les Alpes, région montagneuse, sont particulièrement sensibles au changement climatique, et le cycle de l’eau y est bouleversé. Le manteau neigeux s’amincit aux basses et moyennes altitudes et fond plus tôt au sortir de l’hiver, tandis que les glaciers, eux, deviennent déficitaires, accumulant moins l’hiver qu’ils ne fondent l’été. Cela peut conduire dans un premier temps à des inondations au printemps, et cela soulève à plus long terme des interrogations sur le débit des cours d’eau dépendants de la fonte des glaciers. Par ailleurs, bien qu’il n’y ait pas de consensus à l’heure actuelle, il est probable que l’arc alpin connaisse une augmentation de la variabilité météorologique avec un enneigement plus irrégulier, plus de sècheresses et de canicules, ainsi qu’une augmentation de l’intensité des précipitations extrêmes.

En somme, tous les échelons de la société risquent d’être impactés dans cette région en matière d’accès à l’eau, que ce soit pour la ressource en eau potable, les besoins en agriculture, les usages du secteur de l’énergie – avec le refroidissement des centrales et l’hydraulique –, mais aussi les loisirs, sur lesquels repose actuellement une grande partie de l’économie des zones de montagne.

\vspace*{\fill}



\newpage
\begin{adjustwidth}{2.2cm}{2.2cm}

\begin{center}
\textit{\textbf{Lever de rideau}}
\end{center}
\begin{spacing}{0.8}
~
\end{spacing}
Pourquoi ? Je vous le demande bien tiens, pourquoi ? Questionnement enfantin, mais si spontané après tout, peut-être même instinctif, ou du moins dès lors que l’on accède au langage. Il nous bouscule avant de nous faire vibrer. À chaque « pourquoi ? » c’est une bascule, un pas vers l’éternité. On a parfois l’impression de l’avoir digéré, et pourtant, on le couve toujours en réalité. Il n’est en fait jamais totalement satisfait, et un relent finit tôt ou tard par resurgir. Un « pourquoi ? » est très certainement synonyme de bonne santé mentale et trouve sans cesse matière à se déguiser pour nous suivre, partout. Tantôt partagé, tantôt intime, il sert notre ignorance à toutes les sauces en cachant une délicate subtilité qui, au premier abord, demeure tant imperceptible qu’inaccessible. À chaque apparition, c’est une visite de nos entrailles, une entaille dans nos tripes. Avant qu’il nous inspire, on le soupire ; pour qu’il nous régale, on l’avale. Alors oui, parfois on préfère n’en faire qu’une bouchée, le gober pour y répondre grossièrement. Mais l’expérience nous montre que refoulé, il nous ronge et finit généralement par être régurgité, pas toujours au bon moment, et souvent bien plus acide qu’il ne l’était initialement. Ainsi, même s’il dérange nos papilles, chatouille nos narines, humidifie nos yeux, un « pourquoi ? » reste un os et mérite donc qu’on prenne le temps de mastiquer. On s’y casse les dents pour finalement déguster un pauvre, mais ô combien riche de sens, « parce que ». Déclic de la mise en action, il nous ouvre les portes de la sérénité. D’un passé digéré, c’est la faim de l’avenir qui revient. Lorsque l’on goûte à la substantifique moelle du « pourquoi ? », la compréhension nous ouvre l’appétit des fameux « pour quoi faire ? », « comment ? », « avec qui ? » et « quand ? ». Demain, je me mettrai en marche.
\begin{spacing}{1.2}
~
\end{spacing}
\begin{center}
\textbf{\textit{L’union du départ}}
\end{center}
\begin{spacing}{0.8}
~
\end{spacing}
Tous réunis, là, au même endroit, au même instant. Tout autour, les cimes se dévoilent et les crêtes se dessinent. Tableau noir et blanc d’un instant enivrant, étourdissant. Un nouveau jour se lève, et déjà ce sentiment grisant d’être à l’aube de quelque chose de grand, d’historique. La nuit a été courte pourtant. Des échanges sans fin pour régler les derniers détails, être fin prêts, rivés sur l’objectif, décidés, déterminés, peut-être même déraisonnés. Mais c’est justement cette folie que nous sommes venus chercher ici ; s’imposer ce qui jusqu’alors nous paraissait impossible. Sortir de notre zone de confort comme certains diraient. Après tant d’ébauches, enfin l’action. Après tant d’obstination, parfois par intérêt, parfois sous la contrainte, c’est un but commun que l’on chevauche et qui nous réunit finalement en ce lieu. Mais sommes-nous prêts ? Sommes-nous vraiment crédibles ?

Je suis déjà quelque peu ému de tout ce chemin parcouru pour en arriver là. Impressionné d’être ici, je me laisse saisir par l’instant, par la beauté du cadre et les sourires de la foule. C’est une vraie fourmilière, ça grouille de partout ici-bas, et pourtant, si l’on porte le regard au loin, tout semble immobile, massif, vertical. Magnifique. Cette montagne, échappatoire de mon enfance et maison de mes rêves, était aussi le berceau d’un amour passionnel et fusionnel. Une relation qui dévorait la fougue pour servir les saveurs subtiles du bonheur. Elle donnait frivolité à la fatalité, merveille au néant. Mais c’est aussi au cœur de cette même montagne que les cimes de la vie auront séparé nos chemins. Depuis, j’ai la sensation que tout me ramène à elle, chaque pensée, chaque geste, chaque intention. Sans sa présence, je me sens étrangement vulnérable, mais si ces sommets nous ont séparés, alors je les gravirai, seul s’il le faut, pour la retrouver, là-haut.

En ce matin d’été, alors que le soleil est encore caché derrière les montagnes, je me laisse émerveiller par l’évolution des teintes pastelles du ciel et l’apparition des contrastes de couleur sur les hauteurs. La douceur saisissante de ce qui m’entoure tranche avec l’adrénaline qui s’empare peu à peu de tout mon corps. Elle me pique et me fait frémir.

Nous y sommes. Nous avons eu beau nous projeter, faire, défaire, refaire ces scenarii dans notre tête, l’inconnu et l’incertitude semblent à nouveau frapper. Tant de fois secoués par ces humeurs taquines qui se répondent et se contredisent, de l’espoir à la crainte, de la peur au courage, de l’ambition à l’inhibition, des obstacles aux rêves. À chaque pas en avant, ce sentiment de glisser en arrière. Force est de constater qu’imaginer, conceptualiser, concevoir ne reste qu’un cheminement ancré dans l’abstraction, bien loin du réel. Ce n’est que plan, idéal, utopie jusqu’à cet instant, jusqu’à cette bascule de la pensée à l’action. Longtemps visée, elle est aujourd’hui la seule option.

Alors que notre volonté commune est en passe d’être entérinée, plus question de se sauver, il faut maintenant faire face. Mais tandis que l’on brûle d’envie, les risques nous soufflent et les aléas crépitent ; c’est tout l’édifice de pensées bâti à la hâte pour aller cueillir cet objectif qui semble s’embraser. Cette réunion est somme toute la clef de voûte de cette maçonnerie. Les épaules s’affaissent, les regards s’abaissent. La quiétude devient soudain quête, et l’inquiétude tempête. Les cœurs palpitent, la peur nous paralyse. Terrifiant.

Pourtant, tout est né d’une motivation qui nous a tous saisis à la gorge, ce quelque chose qui nous obnubile, qui nous accompagne jour et nuit. Cela fait tant d’années que nous l’envisageons, tant d’années que nous le préparons. Cet objectif est un monstre. Nous l’avons construit sans jamais accepter de le regarder en face. Nous n’avons jamais eu le courage de le dompter, il nous a toujours repoussé. Notre ambition, jusqu’alors déception, a été fantasme avant de sombrer dans le marasme. Après tant d’échecs, ce nouveau départ doit donc être célébré à la hauteur du défi. Qu’il soit beau, qu’il soit grand, comme ce géant de glace, là-haut, perché, qui a façonné son lit dans la roche des millénaires durant. Qu’il soit tenace, qu’il soit puissant pour sceller le cercueil de nos peurs.

La tension reste malgré tout palpable et les visages malléables. Il y a affluence et ça joue des coudes pour faire sa place, marquer son territoire. Une petite odeur d’hormones vient d’ailleurs me picoter le nez. Ici encore l’atmosphère est très virile, les choses n’ont pas vraiment changé… Difficile de percevoir le vrai du faux dans tout ce que l’on entend subtilement autour en tendant l’oreille. Croire ou ne pas croire les compétiteurs au chapeau melon qui partent plus confiants que ce que leur cœur pastèque ne peut effectivement supporter comme pression. Croire ou ne pas croire les rêveurs à la douceur kiwi qui pensent partir tranquillement en balade et qui risquent de se prendre, en chemin, une pêche en pleine poire. Croire ou ne pas croire les pressés couleurs citron qui finiront très certainement par subir le jus acide de leurs angoisses. Croire ou ne pas croire les charrieurs et leur sourire banane dont l’amertume pamplemousse giclera à la moindre écorchure. Et puis il y a le reste de la corbeille, sans grande prétention, certains pas encore mûrs, d’autres un peu trop. Qui croire ? Peut-être soi ? Dans sa juste mesure si possible, idéaliste mais réaliste, craintif mais primitif. Un mélange explosif de confiance et d’appréhension, cocktail de l’intox. Chercher à prendre la mesure du défi sans mépris, à être lucide sans effroi.

Les discussions sont maintenant closes, le soleil nous illumine et la température grimpe. Le départ est proche et son imminence rend les gens de marbre. Le bras levé de l’organisation déclenche momentanément un vent glacial, et d’un coup sec symbolique, la glace est brisée signant le grand départ. L’excitation nous submerge, l’immersion démarre. La course est lancée dans une certaine confusion où se mêlent des applaudissements et des cris de joie, d’inquiétude aussi à bien des égards. Instant galvanisant où la peur devient rage et le courage, ferveur. Frissons garantis. Telle une vague qui se soulève et qui déferle dans un chaos ordonné, nous partons en route vers l’inconnu bien que la même destination nous réunisse tous. Peu importe notre état, notre préparation, ou encore nos facultés. Plus le temps pour les questions et les projections, nous plongeons maintenant dans l’action, dans le présent qui s’offre à nous, et nous n’avons déjà plus le temps. Car le chemin est jonché de points de passage obligatoires que l’on doit atteindre à temps. Régulièrement espacés, ils scindent le chemin à parcourir en étapes intermédiaires. Sans eux notre détermination vacillerait. Ils nous portent et nous guident dans ce dédale de faiblesses tant l’objectif est loin. Ce sont à chaque fois des réconforts sans égal, des moments forts et frugaux. Dans la générosité, ils nous rechargent ; dans le partage, ils nous motivent. Grâce à eux, nos limitent tombent et le dépassement devient plaisir.
Nous voilà donc lancés dans une course contre la montre, ou plutôt une course avec les autres et contre nous-mêmes. Les yeux rivés au loin, nous nous mettons alors tous en chemin, ensemble. L’envie est là, partout, omniprésente, elle dégouline de bonnes intentions. Avancer coûte que coûte vers l’objectif.
\begin{spacing}{1.2}
~
\end{spacing}
\begin{center}
\textbf{\textit{La soif}}
\end{center}
\begin{spacing}{0.8}
~
\end{spacing}
L’instant collégial du départ laisse instantanément place à un démarrage en fanfare pour certains. Partir vite, trop vite peut-être pour beaucoup. L’émulation collective donne des ailes pour faire ces premiers pas. Le rythme est rapide bien que le chemin soit long, si long... Très vite, le souffle devient court et le cœur s’emballe. La volonté bute déjà sur les capacités, les écarts se creusent. Dès à présent, dépasser est une fierté et se faire dépasser, une cruauté. Il faut déjà bien souvent avancer seul, trouver son allure, faire ses preuves. Je rentre dans ma bulle, tout commence.

C’est à ce moment-là que le plaisir, fondateur, est aussi moteur. Lorsque l’on attache du sens autour du cou d’un défi, il nous tire et nous propulse. Tout est si beau, si grand, si vaste. Je me sens libéré, comme si tout était accordé, parfaitement orchestré, une osmose. Je ressens un bien-être profond. Tout ce que j’entreprends maintenant était pensé, préparé, et me procure un puissant sentiment d’accomplissement. Je m’élève tant j’ai la sensation d’être invincible. Rien ne peut me résister, pas même ces obstacles rencontrés en chemin que je franchis avec hargne pour gravir les sommets. Là-haut je domine, là-haut je survole. Je savoure ces petits plaisirs en apesanteur.

Une fois encore, je suis pris de cette soif ardente d’être en son sein. Elle nourrit mes désirs et m’abreuve de bonheur. C’est une passion indicible, imperceptible qui nous lie, un ressenti, un rêve. Voler, toucher les étoiles. Je suis satellisé, je plane, j’erre dans cet espace sans limite, où ma raison de vivre est ma passion. Je suis fou, fou d’elle, elle est si belle, si harmonieuse, toutes ces couleurs, toutes ces odeurs, tous ces sons. Elle me comble et je succombe à chaque fois. C’est elle finalement l’objet de ma motivation, la montagne.

Crétin des Alpes va ! Toujours prêt pour une envolée lyrique quand la souffrance ne le défie pas. Mais quand il faut faire face, affronter la douce difficulté et accueillir la dure réalité, il n’y a plus personne. Encore dans le refus, comme aveuglé, alors que les paysages, désertiques, sont méconnaissables. À chaque pas, derrière le sol que l’on piétine, un petit nuage de poussière s’élève. Sèche, fissurée, craquelée, la terre elle-même semble s’évaporer. Son squelette de pierres se dévoile, plus rien ne coule. La terre ne respire plus, elle se meurt. Et que dire de ce lac ? À l’époque, il était toujours au plus haut l’été, mais depuis des années maintenant il se vide sans jamais se remplir d’autant au sortir de l’hiver. L’eau s’épuise, la vie se volatilise. Ce lac, aussi artificiel soit-il, a son charme lorsqu’il est à niveau. Il fait vivre certains, irrigue les terres l’été et nous alimente en énergie une bonne partie de l’hiver. Son utilité avait finalement le bon côté de se parfaire d’une image attractive en période estivale. Mais là, plus rien, desséché, il laisse en son lit de mort un paysage aride. 

Amorphe, je repense à ces souvenirs, et je me mets à pâlir devant les conséquences des actes passés, aujourd’hui visibles. De mon inclination servile de partir à ta conquête, ai-je fini par tenter de t’acquérir ? Toi, oui toi, t’ai-je bousculée ? t’ai-je étouffée ? Car ce que tu as refoulé s’avère soudain refaire surface. De mes excès de liberté, tu es tempête. Perdue tu l’es, je le vois, je le sens, tu me le montres. Tous ces changements brutaux t’ont mise hors de toi, hors d’équilibre. Instable tu l’es, imprévisible tu le deviens. Je ne te reconnais plus, je ne te saisis plus, je ne te comprends plus. De ce quelque chose en toi qui te pèse, tu sembles en pâtir et tu parais t’en protéger. Tu me rejettes comme si nos chemins ne s’étaient encore jamais croisés.  Dis-moi… Ce partage sans limite, cette passion, cette fougue, cette complicité ; cette liberté d’être avec toi, cette liberté que tu m’octroies, en ai-je finalement abusé ? Dis-moi… Cette communion sans frontière, cette magie, cette hystérie, cette excitation ; cette générosité de la première heure, à force de puiser en toi, ai-je fini par l’épuiser ?

Le temps passe et la sensation de trépasser se fait grandissante et pesante. Je suis comme plaqué par cette chaleur étouffante, suffocante, accablante. L’air est tout bonnement irrespirable, il m’asphyxie. Plus de place pour le futile énergivore, tout se concentre, tout s’articule pour aller de l’avant. Le superflu se dissipe, ne reste plus que l’utile pour avancer et le beau pour prendre du plaisir. J’ai la sensation de me densifier, de m’ancrer dans l’instant, dans le territoire. Car je dois déjà puiser au plus profond de moi. Entreprendre ce qui au début paraissait dérisoire est dès à présent presque insurmontable. Chaque effort demande dorénavant un investissement total. Vivre, ou peut-être même survivre, désormais. Étancher ma soif n’est plus un plaisir, c’est un besoin auquel répondre, là, maintenant. Plus question de gaspiller l’eau, les réserves sont si pauvres que cela en devient critique. Pourquoi ne pas avoir anticipé ? C’était une évidence qu’elle finirait par manquer… Pourquoi ? C’est maintenant une obsession vitale... Toutes ces années à tourner le robinet sans se rendre compte du trésor que l’eau représente. Toutes ces années à passer sous silence les murmures de ses absences qui finiraient très probablement, un jour, par être des murs à franchir. La réalité ne connaît pas la compassion. Si fautifs, nous sommes, alors elle nous court après et finit toujours par frapper, ce n’est qu’une question de temps. Ici, c’est la chaleur qui assomme et l’air sec qui abat. Tout semble s’être vaporisé, moi le premier… Je ne suis plus que l’ombre de moi-même, fantôme de mes songes. La gorge desséchée, je titube, je me décourage. J’ai soif.
\begin{spacing}{1.2}
~
\end{spacing}
\begin{center}
\textbf{\textit{La faim}}
\end{center}
\begin{spacing}{0.8}
~
\end{spacing}
Dans la difficulté et la persévérance, la vue se floute et les larmes coulent. Les objectifs fondent, et les rêves sont siphonnés. Je m’y accroche quoiqu’il en coûte, mais la machinerie s’enraye, tout se dérègle. C’est l’énergie qui manque maintenant... Pourquoi une telle épreuve ? Pourquoi ? Pourquoi tant de souffrance ? Pourquoi ?  Les émotions me submergent, et ma sensibilité est soudain décuplée. C’est un véritable raz-de-marée émotionnel, une vague ravageuse qui me noie et me laisse sans voix, je bois la tasse. Me voilà maintenant jeté sous une cascade de remords.

Ils avaient pourtant prévenu au dernier point de passage que la suite allait être difficile et qu’il fallait autant que possible faire ses réserves. Ils le disaient déjà bien avant même, mais une fois encore, je ne les ai pas écoutés. J’aurais pu, j’aurais dû… J’aurais pu, j’aurais dû anticiper, prévoir ces cas extrêmes. Je les savais possibles, mais j’attendais peut-être tragiquement ce vil plaisir de les vivre. Perte de lucidité ou excès de zèle ? Le résultat en est malheureusement le même. Plus d’eau, plus de nourriture. Rien de plus vital, rien de plus fatal, car avec la soif, la faim, plus rien ne fait corps. D’une priorité piétinée, le doute se répand. On me passe, me dépasse sans que je ne puisse réagir. Détruit, épuisé, ce ne sont plus les feuilles qui tombent, c’est mon corps qui me lâche, c’est mon âme qui s’effondre. Hécatombe. Plus rien ne me retient, la chute est proche, vertige de l’abîme. Je ne suis plus, je suis absent, je ne suis que néant. Abattu, déchu, accablé par le ciel. J’ai faim.

À vouloir la dompter, elle a fini par avoir le dernier mot. Elle a été crue, franche, sèche. Elle m’a abandonné, jeté à terre, et me fait maintenant transpirer, baver. Transpercé en plein cœur, le goût du sang inonde ma bouche. Je suis déchiré, à cœur ouvert. C’est l’effondrement brutal d’un rêve, d’un idéal que je m’étais construit, soudain inaccessible. C’est une porte qui se referme, une porte qui claque même. Elle me claque à la face. À force de mépriser mon être, de salir l’essence de la vie, un point de non-retour m’attendait. À force de m’y méprendre, il m’a crûment surpris. Ma confiance et mon orgueil sont pressés avec un tel entrain qu’il n’en reste qu’une pauvre bouillie. Comment sortir la tête de l’eau ? De l’air s’il vous plaît, je suffoque…
Subitement insensible, inerte, ou presque, le milieu dans lequel j’évolue semble m’avaler. Écrasé, oppressé, mon être se liquéfie, je me dissous. Autour, tout me paraît mou. Le paysage a l’air de flotter, et la ligne d’horizon elle-même semble flasque. Plus rien de ce qui m’entoure ne m’affecte, je suis là sans être, je suis ailleurs, nulle part. Je fais face à un vide abyssal, je le sens. Creux je le suis, j’ai perdu ma substance. Pire, mon esprit semble se dissocier de mon corps. Je ne ressens plus que le poids de mon âme, si lourde, si pesante… Est-ce le poids de mon passé ? J’ai en tout cas l’étrange et désagréable sensation d’être en rupture totale avec ce qui m’entoure, avec moi-même aussi. Reflet douloureux, radical, bouleversant d’un regret qui laisse des traces.

Entre un avant où l’on pense et un après où l’on agit, il n’y a qu’un pas, mais il est si compliqué à entreprendre… Je peine à être les jambes de ma pensée, à fournir l’énergie de mes souhaits. Je finis même par me perdre. Où suis-je ?

Face à la difficulté, le sens dégouline malgré tous les efforts jusqu’alors consentis. De l’objectif initial je me lasse, l’adversité m’enlace. Je n’ai plus la force. Penser m’est impossible, seule la discipline me permet d’avancer malgré tout. Un élan de survie. Plus rien ne me traverse, plus rien ne s’exprime. Je suis pris à la gorge par l’incompréhension, noué par l’injustice.  Le sens devient soudain démence, un tout qui se divise… Ô toi… moi… je ne sais plus… Sur moi, tu fais subir les affres de ce qui t’enchaîne. Sur moi tu te déchaînes. J’essaie de saisir les frustrations que tu exprimes, j’écoute les libertés que tu revendiques. Mais au fond, je souffre de voir que cette passion raisonnée s’éteint malgré moi. Détresse d’une tornade intérieure. Plus rien ne se pose, pas même l’esprit. Oui… tu sais ô combien je pense à toi, même dans la rupture… Tu sais à quel point tu me fascines, même dans l’incompréhension… Tu sais si bien comme j’ai envie de partager avec toi, même dans la douleur… Tu sais comme nous nous correspondons, même dans l’éloignement… Je ne sais plus qui je suis, je ne sais plus si tu es ; est-ce que nous sommes ? Sans réponse…

Je suis blessé, profondément, prêt à tout plaquer, mettre à l’eau tout ce que j’ai déjà accompli, tout. Je n’en peux plus, c’est fini. J’abandonne. Trop dur, trop long, trop… c’est trop. Ça suffit, je suis à bout. Tout me fait souffrir, tout m’écrase, tout m’absorbe. Je m’efface. Mon mental se fracture. Pourquoi moi ? Pourquoi ? Je suis frêle, fragile, faible, vulnérable. Genou à terre. Comment garder son sang-froid quand tout bouillonne ? Comment nourrir cette larme d’espoir, ce chagrin de victoire ? Je bute… Lutter je ne peux plus… Mais non, non et encore non ! Je me dois d’assumer, d’admettre mes torts, d’endurer l’effort. Il y a des responsabilités que l’on tarde à endosser, certaines que l’on accepte dans l’échec. Je ne peux pas, je ne veux pas abandonner.

La persévérance finit d’ailleurs par être récompensée par une rencontre inattendue en chemin, une rencontre qui redessine le destin. Un regard suffit à lui faire comprendre ma misère et mes besoins devenus vitaux. En retour, dans un élan de générosité salvateur, on m’offre enfin ces formidables gorgées d’eau tant désirées, ces extraordinaires bouchées d’énergie tant espérées. De cet échange déséquilibré et impromptu, j’en tire une véritable bouffée d’air frais. J’en garderai la mémoire, car ce don est un sauvetage dont rien en retour ne saura être à sa juste valeur. C’est un geste qui abreuve l’égo d’un jus d’amitié, d’un zeste d’humilité. À toi qui m’as doublé, à toi qui m’as relevé, merci. 

En trouvant satiété, j’ai perdu ma cécité. D’un premier pas différent des autres, d’un premier pas vers l’autre, j’ai rompu avec mes excès pour m’attacher aux plaisirs les plus simples qui soient. Un rien rend heureux, et un sourire donne tout. Ce n’est peut-être qu’un état d’esprit, qu’une façon de voir la vie, les autres, de s’y lier sans s’y emmêler par intérêt. Euphorie d’un instant, extase d’une rencontre. Je reprends goût, je reprends vie. Je réalise que l’on se délecte du talent, mais que l’on jouit finalement du courage. J’ai hésité à baisser les bras, je me suis obligé à revoir mes objectifs, mais je ne lâcherai rien, je donnerai tout. N’est-ce donc pas par amour de l’autre, donc de soi par effet miroir, que l’on avance, que l’on déplace des montagnes ?
\begin{spacing}{1.2}
~
\end{spacing}
\begin{center}
\textbf{\textit{Le délire}}
\end{center}
\begin{spacing}{0.8}
~
\end{spacing}
Tous ces tracas derrière moi, me voilà donc reparti sous les dernières lueurs du crépuscule. Abandonner ou continuer ; poursuivre ou fuir. L’incertitude de cette décision a laissé place à la solitude d’une rédemption qui me plombe, car je reste malgré tout en état de choc. Ma flamme se consume doucement, et l’iceberg de mes doutes me glace le sang. De la partie émergée que sont mes peines, je cache en réalité un véritable cancer, immergé, invisible, profond. Je couve un tourment discret, dont les pensées noires n’ont aucun mal à me renverser comme pour me défier. Alors, qu’est-ce que tu vaux ? Alors, tu te croyais plus fort ? Chaque pensée est une fessée, à l’image de la fatigue et de l’obscurité qui s’abattent désormais sur moi.

Je me laisse maintenant guider par ce faisceau de lumière, lueur d’espoir qui, tel un guide, pave le chemin vers l’objectif. Un tunnel dans la nuit, dont plus rien ne sort, pas même mes larmes jusqu’alors. Je suis tétanisé, mais je m’accroche aux vestiges qu’il reste du plan initial, complètement soufflé par l’avalanche de douleurs rencontrées depuis le grand départ. Heureusement, face à ces difficultés, mon indifférence givre ma perte de sens, et ma discipline gèle mon spleen. Mais mon corps me joue toujours des tours, symptomatique d’une volonté qu’il ne parvient plus à satisfaire. À chaque pas, c’est une porte que l’on force en moi. À chaque pas, relancer, encore et encore. Toujours relancer. Toujours ce même plan qui se tient en un mot… relancer. Encore relancer. Je me dépasse à en perdre littéralement la raison. Je débranche mon cerveau.

La machine avance mais l’esprit dort. Bringuebalé par les vents, il vagabonde, et l’arrivée sur les cimes le fouette tout à coup. C’est une véritable gifle glacée. Un réveil ? Je ne sais dire, mais il vacille tant cela paraît palpable. Les yeux se sèchent un court instant puis se brouillent, et soudain, au loin, quelque chose d’un blanc éclatant. Dans ce clair-obscur illuminé par la lune et les étoiles, c’est un monstre de glace qui me fait face. Perception étrange de voir en face le froid soudain saisissant que mon corps ressent. Plus j’avance vers lui, plus il semble s’éloigner de moi. Irréel. C’est une véritable majesté onirique qui se dévoile. Il s’en dégage une telle puissance, une telle sérénité. Là-bas, au loin, si imposant, si majestueux, ce glacier m’appelle. Je me rappelle alors mes premiers pas sur cette meringue gelée à contempler ses formes, à écouter ses sons, à relever ses imperfections qui le rendent merveilleux. Et dire qu’il m’a fallu toutes ces années avant d’apprendre que cette masse n’était pas immobile, qu’elle s’écoulait, qu’elle répondait en fait à son environnement, et qu’un glacier pouvait donc se révéler fragile.

Quelques instants plus tard, je me frotte les yeux humides, et plus rien. Glaçant… Ce géant de glace, là, devant moi, qui m’éclairait soudain de sa réflexion, n’est plus, plus là. J’hallucine… Comme un rêve, ou plutôt comme un souvenir qui refait surface, celui d’une époque que j’ai connu mais désormais révolue. Je sors brutalement de ma torpeur, car on m’avait dit que tu étais éternel, au moins à l’échelle d’une vie humaine, et voilà que tu n’es plus. M’a-t-on menti ? Ou avons-nous fini par faire mentir ce que nos parents et grands-parents contaient. Il n’y a plus l’ombre d’un doute je crois… J’en suis même certain, et il m’est aujourd’hui insupportable de me résigner en tant que témoin répréhensible plutôt que de me désigner comme acteur repenti. À moi d’accepter de recevoir le témoin d’une génération à contre-temps pour prendre le relais d’une course qui se joue à plein temps dès à présent. Chère glace, peux-tu rester, pour toujours, la compagne de mes épopées idylliques plutôt que de faire partie, à jamais, de récits mélancoliques ?
Je suis en colère, fou de rage. Quand d’un souvenir, le réveil sonne coupable. Elle fait mal cette responsabilité. D’une culpabilité que je porte, je m’emporte. Tout ça pour ça ? J’ai mal. Je ne parviens plus à retenir ces relents de « pourquoi ? » qui ne trouveront jamais de réponse. Cette disparition matérialise la mienne. Piqué au vif, j’ai le sentiment d’avoir été, un temps, privilégié, et d’être, maintenant, réfugié d’une ère à présent achevée. Je crache l’incompréhension, je déglutis l’injustice. Après avoir tout tenté, tout donné, tout laissé dans la bataille, je dégueule. D’un rêve déçu que mes enfants, un jour, voient ça aussi de leurs propres yeux, chez nous, là-haut, je murmure mes regrets. D’un passé déchu qui devient un supplice, je crie ma peine. 

Je ravale mes idées noires, et je me laisse bercé par le ciel étoilé. Il reste sûrement ce qui nous lie, toi et moi, encore et toujours. Quand mon regard s’y perd, c’est ta présence qui me gagne, c’est une sérénité que l’on partage. Discrètement je t’admire, délicatement tu nourris mes rêves. Une étoile filante passe. Et si cette chance d’avoir vécu ce temps-là était une invitation à s’engager pour demain ? Et si cette reconnaissance d’avoir profité dans le passé était une promesse à faire pour l’avenir ? Et si finalement, d’un texte ratifié, une nouvelle histoire pouvait s’écrire, pour l’éternité ?
Dans l’obscurité, les points de passage se succèdent. Ils s’illuminent à chaque fois d’échanges spontanés, vrais, francs. Ils ont beau être éphémères, ils sont un véritable point d’appui, un réel soutien dans l’effort entrepris. Ils permettent de ne pas se désunir, de tenir. Sans celles et ceux qui en sont les garants, l’objectif ne pourrait être atteint, il se briserait en mille morceaux. Alors à chaque fois que j’arrive à l’un d’entre eux, je prends mon temps. J’y consomme mon amertume, soigne mon hémorragie et réapprends à croquer la vie. Quelle que soit alors ma situation, j’y retrouve toujours ce quelque chose qui fait chaud au cœur, qui m’anime, m’enivre et me rend finalement heureux. 

Mais dès que je repars, dès que je m’éloigne, tout semble à nouveau s’éteindre si ce n'est la flamme allumée au départ. Dans le dur, elle perdure, dans le noir, elle reste espoir. Et pourtant je me lamente, toujours, je désespère, encore. J’ai le regard vitreux, triste, hypnotique. J’avance maintenant enchaîné au passé, tout en m’efforçant de braver la douleur pour cicatriser. Conscient mais anéanti par cette situation qui me dépasse, j’ai la sensation de ne plus tenir les rênes, de toujours avoir un temps de retard, de subir les dommages des joies d’autrefois. Je me complais finalement dans un état d’esprit en demi-teinte, ni blanc, ni noir, sans couleur peut-être, rien de tranchant pour sûr. C’est une forme de fatalité qui me gagne et que je bois naïvement comme un remède. Chaque pas est un souvenir qui se rallume. Chaque pas est une brûlure que l’on ravive. Il faut renoncer, accepter de dire adieu à tous ces moments que je ne revivrai plus, devenus frustrations. Crucifier mon passé pour ne pas le refouler, ni l’oublier, mais geler mes larmes que le temps finira peut-être par sublimer. Contempler ces vestiges qui m’avalent pour faire en sorte que ma colère ne devienne ni haine, ni vengeance, que mes souvenirs cristallisent pour ne plus être des coups de poignard. Faire le deuil de mes rêves. Marcher ou guérir, marcher pour guérir. Finis les cris de rage, les manifestations de dépit, le mouvement se réordonne, le calme refait peu à peu surface. Étourdi et redevenu insouciant, tout semble maintenant figé, à commencer par ma mémoire, comme momifiée.

La nuit a éteint le feu et voilà que le petit matin le rallume. Bascule des ténèbres à la lumière ; déconstruction d’une croyance pour repartir de rien, repartir d’un rien, et en reconstruire une nouvelle. L’aube est un nouveau départ. Il est peut-être là le virage. Changer de cap. D’un passé devenu répulsif, j’ai placé mon espoir en l’avenir, dorénavant attractif. Ronronnement d’une rancœur qui nous perd, réveil d’une rage nouvelle qui nous guette, celle d’aller cueillir demain. Le moment de l’acceptation est arrivé, laissant derrière lui une période en ruines, glaçons que le temps, lui seul, fera fondre. Je reprends enfin mon souffle, je l’écoute, il me rythme. Le soleil m’accueille, il me caresse la peau et me redonne de l’énergie. Il m’offre cette affection dont j’avais tant besoin, cette forme de quiétude que je recherchais. Je saisis alors la subtilité et la pudeur de notre relation qui nous lie l’un à l’autre, aussi imperceptible soit-elle. Ce n’est que flux d’énergie, un équilibre, ou pas, lorsque le système est en marche forcée.
\begin{spacing}{1.2}
~
\end{spacing}
\begin{center}
\textbf{\textit{La paix}}
\end{center}
\begin{spacing}{0.8}
~
\end{spacing}
Depuis quelques temps, j’ai l’impression de moins subir, mais il me semble que j’encaisse sans réellement affronter l’âpreté de la situation. Mon sort est entre ses mains, alors j’accepte la douleur de l’incertitude, et me satisfais d’une condition quelque peu désespérée. Mais son manque me terrasse, m’assaille de coups dans le ventre, je tressaille. Il broie ma chair, tord mes boyaux et finit par saucissonner ma confiance. Je suis littéralement séché pour rebondir. Est-ce que je ne le peux pas ? Ou est-ce que je ne le veux finalement pas ? Peut-être suis-je libre de vouloir, ou peut-être suis-je captif de devoir. Elle me manque tant… C’est elle que je veux ! Alors jusqu’à quand pourrais-je maintenir une forme de résilience face à une telle absence bien qu’elle me soit vitale ? Cela ne peut qu’être temporaire, je le sais. Donc je doute, oui je doute de la revoir, la pluie, car mon esprit reste, malgré moi, fuyant, incapable de se poser. La nostalgie comme thérapie, nouvel état d’âme d’une apathie qui se blâme. Anesthésié par l’effort, hypnotisé par les souvenirs, j’avance tête baissée, le regard dans le vide. Mes songes s’ancrent de nouveau dans le passé, mes pas dans le présent, un bourbier pour l’avenir qui couve néanmoins une énergie créative nouvelle, où la douleur inspire et la peine s’exprime.

Plus rien ne m’importe si ce n’est moi. Je recentre et me concentre pour mieux me retrouver. Je recherche la beauté du geste, la beauté de l’être, rien d’autre, simplement être là, et vivre ça. En me réappropriant mon imaginaire et laissant exprimer ma créativité étonnamment débordante, j’accède à un bien-être profond. Je me sens bien, léger, souple, aérien. La forêt qui m’entoure me captive. En fait, je réalise que cet objectif, de par le sens qu’il porte, est devenu pour moi un devoir, et qu’il parvient tout de même à se parfaire d’une liberté. Je me sens libre, libre de mes choix, libre de mes actes. Je sens que l’audace se réveille en moi.

C’est une odeur quelque peu dérangeante, gênante, cuisante même, qui me pousse soudainement à sortir de ma rêverie mélancolique, de mon hibernation. Une odeur synonyme de danger. Le brûlé. Oui, c’est ça, ça sent le brûlé ! L’odeur est saisissante, ça sent la mort ! Dans un instinct de survie, je relève les yeux face à, je le découvre, ce qu’il reste d’un incendie dévastateur. Se révèle à moi un paysage calciné, même si la vie, imperceptible à première vue, subsiste peut-être encore sous certaines formes. Les souvenirs douloureux semblent être devenus cendres, brûlant un temps, ils sont devenus volatiles après avoir été consumés par le temps. Ne restent alors que certaines fondations, certains troncs à nu, encore debouts, sans écorce. Quel choc de voir que cela… s’étend à perte de vue, jusqu’en haut des cimes. Et… Non, impossible… Comment ? Non… Je reste brutalement bouche bée, totalement stupéfait. L’imprévu frappe à nouveau, je suis complètement pris au dépourvu.

Droit devant, en face, là-bas, mais presque déjà ici finalement, c’est… Oui c’est une rencontre qui se profile, des retrouvailles peut-être hostiles, peut-être fertiles, assurément électriques. Elle est donc de retour, elle, oui elle… Tant attendue, la revoilà… Tant espérée, elle est là… Elle était allée voir ailleurs, mais elle finirait par revenir, je voulais y croire. Tant d’attente pour qu’elle se révèle finalement à moi sous sa forme la plus tumultueuse. Elle semble presque menaçante à se dresser si haut dans le ciel qu’elle pourrait en éteindre la lumière. J’ai comme la sensation qu’elle revient avec une énergie nouvelle, que de mes excès, elle veut m’en servir l’amende. Je ne sais qu’attendre d’elle, mais je la veux, elle m’a tant manqué… 

Après tout ce temps de séparation, j’appréhende de te retrouver, de te redécouvrir somme toute, bien que je te désire, oui toi, au plus profond de moi. C’est comme si c’était notre première rencontre. J’ai l’impression de ne plus te connaître et pourtant le souvenir de ta présence, si proche encore, ravive la mémoire de tous mes sens.  Je te veux, mais je suis saisi d’une angoisse certaine de ne plus être à la hauteur pour t’accueillir. J’ai l’impression de ne plus savoir comment me comporter avec toi. J’ai envie de courir à toi, de sauter de joie, de pleurer de joie, mais ton apparence capricieuse freine mon élan, je me sens presque craintif. Tant d’excitation, mais en même temps de peurs et d’hésitations. Ô toi là, je ne tiens plus, inonde-moi, submerge-moi ! Braouuum... La pluie est de retour, là, droit devant. Elle fête sa réapparition en fanfare, elle déchire le ciel, elle l’électrise, l’orage gronde. 

Elle a changé, je le perçois, elle s’affirme, je le ressens. Mais au fond de moi, je sens quelque chose qui me brûle à nouveau, je sens que de la friction passée peut renaître la chaleur, que d’une étincelle peut reprendre un feu intérieur. Elle a ce quelque chose qui me laisse pantois, toujours, encore, décidément, c’est bien d’elle dont j’ai besoin dans la vie. Mon cœur bat la chamade. Mais elle prend place et semble déterminée, savoir ce qu’elle veut. Elle est autant déconcertante qu’élégante, autant foudroyante que généreuse. Sa volonté suinte, ses rêves perlent. Elle semble s’exprimer à verse, a-t-elle un message à communiquer ? Braouuum… Le ciel gronde, elle se rapproche.

Giflé par de puissantes rafales, je titube. Douché par une pluie torrentielle, je glisse. Et… Scraaaatch ! La foudre, là juste à côté, juste là. J’ai le souffle coupé. Dans un éclair de lucidité, je me jette à terre et me mets en boule. Pourquoi moi ? Suis-je coupable ? Tous ces éclairs, si puissants, c’est glaçant. C’est une pluie de bombes. À quand mon tour ? Le sort s’acharne avec hargne, il m’envoie un signal. La température chute brutalement, et l’eau qui ruisselle sur moi me frigorifie. Le sol desséché, craquelé, est incapable d’éponger cette eau tant attendue. Mon chemin devient en quelques minutes un torrent. C’est une véritable douche froide qui vient éteindre le feu une bonne fois pour toute, lessiver le champ de guerre incendié, lui ôter son odeur de brûlé et mettre en lumière les fondations qu’il reste. Ça tombe dru sur un sol à nu, ça ruisselle sur mon être tout nu. Submergé littéralement par l’instant, sidéré par l’intensité de son retour.

Puis une claque. Un caillou sur ma tête ? Non, une pluie de grêlons d’une violence inouïe. Cette fois elle toque vraiment à ma porte, elle s’adresse à moi, elle me dresse, mais moi je m’accroupis en réponse, terrifié. Elle est déchaînée et me voilà enchaîné à ma peur, pétrifié par cette eau, par ce froid glaçant qui contrastent tant avec la chaleur étouffante et la sècheresse suffocante qui précédaient cet orage. Je ne peux plus bouger. Si vulnérable, mon égo est humilié, rincé, délavé. Mais dans un sursaut d’humilité, je l’écoute, je l’entends. Je ne suis plus là pour lui faire face, mais l’accueillir et m’adapter. À chaque impact, c’est une peine que j’enterre, c’est ma tête qui se redresse. Et les choses tournent même à l’absurde. Le ciel se charge d’un voile blanc cotonneux venant étouffer la tension de ces retrouvailles et apporter légèreté et lumière. Il neige.

J’avais l’impression d’être en guerre sans armes… Mais en guerre contre qui ? Contre moi-même peut-être ? Est-ce le reflet d’une guerre intérieure, une forme de culpabilité que j’alimente ? Elle n’était pas là pour me faire du mal, simplement pour tirer la sonnette d’alarme… Incontestablement, la guerre est une pensée défaite, une perte, tandis que la paix, elle, est une plaie pansée, une fête. Cette neige soudaine, inattendue en est le symbole, c’est la paix qui tombe. Elle me subjugue et m’enchante. Elle est en fait pleine d’attentions et moi d’intentions. En me caressant le visage, elle rouvre mon cœur. Orchestrer un tel retour me laisse sans voix. Car devant cette neige, c’est une fascination qui me reprend, celle d’un amoureux qui n’a d’yeux que pour elle. C’est un regard d’enfant que je retrouve, un émerveillement sans limite. Celui que j’avais à l’époque lorsque le neige tombait, et qui me permettait de rester éveillé, ébahi à la fenêtre, des heures durant pour la contempler, juste pour l’admirer sous les lueurs du lampadaire voisin. Sublime. J’observais alors ce nuage de plumes tomber du ciel et le sol se revêtir peu à peu d’un duvet blanc, étincelant, mettant tout en sourdine. J’étais ébloui par sa beauté, littéralement envouté par son allure désinvolte, flottant et chancelant dans le ciel. Elle ne me demandait rien en retour mais je donnerais aujourd’hui tout pour la retrouver. Chaque année désormais, l’hiver, je l’attends chez mes parents, mais elle se fait malheureusement de plus en plus rare, à mon grand désarroi. Cette neige qui tombe en plein été, aussi absurde cela soit-il au premier abord, est un véritable cadeau, véritable madeleine de Proust que je savoure. J’adore ! Bien qu’éphémère, cet instant de légèreté sous la neige avant que la pluie ne reprenne l’avantage, contient une déclaration d’une beauté sans égal, le suprême, le sublime, le sensationnel d’une situation jusqu’alors contrainte, douloureuse et blessante, le pardon.

À voir l’eau ruisseler, c’est mon sang qui coule à nouveau, je reprends vie. J’ai compris, c’est bon. Ce lien qui nous unit tous, nous l’avons saccagé, pillé, infesté, et pourtant il subsiste malgré tout. Un fil, une ligne, une toile, la vie. C’est un tout qui nous prend à notre échelle, si petite, si infime, si dérisoire. Aussi inspirante soit-elle, la nature a toutes les réponses aux questions que l’on se pose sans jamais nous les offrir. Elle trouve d’ailleurs toujours à s’adapter, non sans forcément nous considérer, bien malgré nous peut-être. Elle est généreuse, car elle ressource sans jamais juger et nous offre toujours sa beauté à portée de sens. Elle se laisse finalement toujours désirer en restant, à jamais, inaccessible. À chaque regard, sur elle, c’est un rêve qui se dépose.

Ô tu me fais bégayer… Car tu es peut-être finalement celle… Oui tu es celle… Celle… Celle qui me fait toucher une forme d’amour tant universel qu’éternel dans sa forme la plus pure, indicible, silencieuse. Je l’ai toujours cherché ailleurs, j’ai toujours cru en son potentiel, mais je n’ai jusqu’alors trouvé qu’illusion. C’est toi que je veux ! Oui, toi qui cherches toujours à t’équilibrer, je t’ai quelque peu perturbée, contrainte, forcée. Mais j’ai compris que tu ne m’appartenais pas, alors assurément, c’est un respect que je t’accorde, et spontanément, c’est un pardon que je t’échange. Je le reconnais, pendant longtemps j’ai abusé de toi, mais j’ai compris que te vivre, ce n’était pas te consommer. Pardonne-moi, je t’en supplie, accepte-moi à nouveau, je te prie. Je souhaite simplement vivre avec toi, en communion. Je t’estimerai, je t’honorerai, alors redonne-moi une chance, s’il te plaît.

Les blessures deviennent cicatrices, car le pardon se cultive pour finalement en cueillir le fruit. C’est une graine d’espoir, commune et apaisée, subtilement déposée et qui attend patiemment d’être abreuvée de nos peines pour germer. Les cendres de l’incendie comme un terreau fertile, rendu fécond par la pluie. Ainsi, des larmes se rallume alors la flamme, la braise se ranime. Brûler nos peines pour arroser notre feu. C’est là un nouveau départ, une nouvelle ère que l’on promeut, une histoire qui redémarre et qui s’exclame. Repartir, différents mais attentionnés. Repartir, complémentaires mais uniques. Repartir, libres mais engagés. Faire en sorte que cet amour soit fécond. Goûtons la sagesse, croquons la simplicité. Trinquons ensemble à la vie pour de nouveau y croire et boire l’espoir. Le pardon résonne comme un éclair, c’est un coup de foudre qui nous illumine, un électrochoc qui permet de retrouver le contact, l’échange. Il nous relie à nouveau, toi et moi, ciel et terre. Reprendre racines pour fleurir à nouveau.
\begin{spacing}{1.2}
~
\end{spacing}
\begin{center}
\textbf{\textit{La solitude de l’arrivée}}
\end{center}
\begin{spacing}{0.8}
~
\end{spacing}
L’objectif, revu à la baisse en chemin, est sur le point d’être atteint. Rarement haï, parfois apprécié, souvent inaccessible, mais toujours désiré, il m’aura mis à l’épreuve et j’aurai fini par faire mes preuves. Tout se fait maintenant machinalement. Je suis vidé, mais pas écœuré, loin de là. Le temps semble d’ailleurs se dilater, faisant sauter une dernière fois les verrous du mental, ou pas. Car étrangement, c’est à ce moment même que profiter prend tout son sens. Je me sens comme un enfant, comme celui que je suis en réalité, un nouveau-né dans ce monde que je perçois d’une manière différente, inédite. Je chante, je ris, je saute dans les flaques et cherche même les plus grandes, car je veux que ça éclabousse mon dos, que ça dégouline sur mes jambes. Je veux être sale, cru, authentique. Cette fougue, cette énergie nouvelle, je la sens, je la prends, je l’exprime, je veux la partager, l’offrir. Plus rien à anticiper, plus aucune tension, plus aucune pression, tout se relâche. Le défi sera relevé, et l’effort mènera à destination. Incroyable ! Qui l’eut cru après tant de péripéties ? Savourer ces derniers moments en chemin, vivre, vivre l’instant, enfin, est la plus belle des récompenses.

Quand la dernière ligne droite se présente, c’est un sentiment de nostalgie même qui m’accompagne. Je donnerais tant finalement pour que cette épopée n’arrive jamais à son terme, faire durer éternellement cet intime sentiment d’accomplissement, tellement mérité. Car le chemin aura été en dents de scie érigeant la montagne comme juge de paix. Il aura fallu que je descende les pans de l’hubris pour gravir les monts de l’humilité. Rien n’aura été facile, mais tout aura été si intense… Avancer, encore et toujours, aura été un combat de chaque instant, entrecoupé de moments de lévitation, de plénitude. Tant de souvenirs concentrés en si peu de temps, tout se bouscule. Cette aventure est une histoire personnelle qui ne se partage pas, ou peut-être que si finalement… Car je ne sais pas si je peux parler d’épilogue tant il y a à raconter sur ce voyage initiatique.

À l’arrivée, les bras se lèvent et le sourire s’étire, mais l’objectif se perd, et une responsabilité se gagne, celle de tirer profit de cette expérience, si riche. L’eau qui coule à nouveau dans les ruisseaux environnants est le reflet d’un changement radical de modèle dont les nouvelles lignes se dessinent. Parti en combattant, j’arrive en pacifiste. Un océan semble me séparer de celui que j’étais auparavant. Car au cours de ma traversée, à la fois si courte et si longue, vingt-quatre heures seulement, mais vingt-quatre heures tout de même, j’ai compris. Oui, en chemin j’ai compris que la compétition n’était pas une fin, mais la faim de l’égo, et que le partage était en revanche adage. J’ai compris que nos connaissances n’auraient jamais la mesure de notre infinie ignorance, et que la souffrance n’était pas une fatalité, mais gage de créativité. J’ai compris que se projeter dans un objectif, ce n’était pas le réaliser, mais que l’atteindre demandait de se réaliser. J’ai compris que conquérir c’était tuer, mais qu’écouter c’était cueillir la vertu. La liste pourrait continuer sans fin, car j’ai compris, in fine, qu’à tes côtés, toi la montagne, tout était à apprendre, peut-être réapprendre devrions-nous même dire dorénavant. À moi d’y croire, il ne sera jamais trop tard. J’ai enfin saisi pourquoi… Parce qu’il y a urgence.

Alors pourquoi se plaindre quand on peut encore enfreindre ? Pourquoi se morfondre quand il est encore temps d’y répondre ? Pourquoi renoncer quand on peut commencer dès maintenant ? Pourquoi se taire quand on peut défaire puis refaire ? Cette liberté de choisir, cette expérience me l’a permise. Je ne suis plus attaché à mes désirs, cloué par mes frustrations, martelé par mes regrets. Je ne suis ni optimiste, ni pessimiste, simplement réaliste en disant que sans elle, oui, sans elle, je ne suis plus. Alors, j’ai fait mon choix, c’est elle que je veux ! Elle est mesure, simplicité, justesse, lenteur, sagesse… La sobriété. Délogée fut un temps, j’aimerais aujourd’hui en faire l’éloge fortuit, car le bonheur s’y loge et rien n’y déroge, j’en suis désormais convaincu. Toi et moi, pour demain, nous ne ferons plus qu’un.
\begin{spacing}{1.2}
~
\end{spacing}
\begin{center}
\textbf{\textit{Baisser de rideau}}
\end{center}
\begin{spacing}{0.8}
~
\end{spacing}
Pourquoi un ultra-trail, encore un ? Encore une de ces courses d’endurance extrême qui nous met à nu. En épluchant nos capacités et disséquant notre volonté, elle nous révèle notre personnalité. Dans un mélange, tant d’espoirs et de gloires que de craintes et de plaintes, c’est un changement de paradigme qui s’opère. Sous couvert de fatigue, de douleurs, elle nous attache à l’essentiel, à l’utile et au beau. Elle illumine nos valeurs et les scelle de principes. De l’adversité éclot la résilience, une renaissance que l’on couve sans jamais jusqu’alors lui donner sa chance. Libération.

Alors bien que cela reste un jeu où le courage est facile et la compétition futile, les erreurs sans influence et l’abandon sans conséquence, faisons résonner le message qu’il porte. Les leçons que l’on en tire sont puissantes et sonnent comme allégoriques de ce que nous devrions faire, à bien des égards, face à la pression anthropique exercée sur le climat et le vivant, face à notre boulimie énergétique, consommatrice et expansionniste. À cela près que la compétition cède à la solidarité, que le courage doit être étendard, les erreurs, une peine, et l’abandon, un point final. Mais à quoi bon expérimenter une telle parenthèse sans avoir le bon sens d’en tirer profit ? La démesure se mesure ici à une forme de fatalité, aux limites des capacités, de la volonté. C’est une leçon de vie, une nouvelle, qui nous ouvre les portes pour bifurquer ; encore considérer nos objectifs comme un mirage, ou bien au contraire, entamer ce virage. Car à quoi cela sert-il d’apprendre si c’est finalement pour délaisser, ignorer, renier le message dans les faits accomplis ? Et si donc, pour une fois, nous refusions la discorde intérieure, la dissonance cognitive ? Et si nous nous accordions et nous appliquions ce que nous avions appris de manière transversale ? Et si au moins nous faisions l’effort d’essayer ? Donnons donc un second souffle à cette parenthèse en apesanteur face au poids de nos responsabilités.

Heureusement, demain le jour se lèvera et marquera un nouveau départ. Seuls, nous tirons les leçons et nous apprenons à nous connaître, mais ensemble nous pouvons passer à l’action et appliquer ce que nous savons. Se présente à nous l’opportunité unique de construire le futur dont nous avons envie. À nous de décider, alors, ce que nous voulons faire. Être frileux pour un avenir capricieux, ou être furieux pour un avenir radieux ? Surmonter des difficultés nous en sommes capables, en y trouvant même un bien-être aussi profond que subtile et suave. À nous donc de nous saisir des enjeux sur le climat en vue d’une réduction drastique de nos émissions de nos gaz à effet de serre, notamment par la neutralité carbone. À nous donc de répondre à l’effondrement du vivant à travers la préservation de sa diversité et de son abondance. Aujourd’hui, renoncer n’est pas une contrainte, c’est une délivrance, une porte qui s’ouvre où les chaînes de nos convictions deviennent actions, et le compromis, dernier rempart.

Pourquoi continuer comme hier ? Pourquoi être à réaction plutôt qu’à propulsion ? Pourquoi attendre de subir quand nous pouvons, encore, choisir ? Pourquoi ? Parce que… Après cette expérience, on comprend que certains « pourquoi ? » n’ont pas de moelle. On a beau les sucer jusqu’au bout, il n’y a rien à en tirer. C’est un pathos qui nous pénètre ; une ataraxie que l’on circlut. C’est une invitation à castrer notre patrie archaïque pour materner un avenir fécond. Parce qu’on comprend bien qu’une introspection est inévitable… Parce qu’on comprend bien qu’il va falloir passer de l’intention à l’action… Puisqu’il faut le dire ainsi, une bascule est nécessaire. Pour quoi faire ? comment ? avec qui ? et quand ? Demain, nous nous mettrons en marche.
\begin{spacing}{1.2}
~
\end{spacing}
\begin{center}
\noindent \textbf{\textit{Bascule}}
\end{center}
\begin{spacing}{0.8}
~
\end{spacing}
\noindent S’imposer ce qui nous paraissait impossible.
\newline Mais sommes-nous prêt.e.s ? Sommes-nous vraiment crédibles ?
\newline Cette bascule de la pensée à l’action,
\newline \noindent Longtemps visée, est aujourd’hui la seule option.
\newline ~
\begin{spacing}{0.8}
~
\end{spacing}
\noindent De notre inclination servile de conquête,
\newline Avons-nous fini par tenter de t’acquérir ?
\newline De nos excès de liberté, tu es tempête,
\newline Instable, imprévisible, tu sembles en pâtir.
\begin{spacing}{0.8}
~
\end{spacing}
\noindent Plus de place pour le futile énergivore,
\newline Les objectifs fondent, et les rêves siphonnés.
\newline Nous aurions pu, nous aurions dû anticiper,
\newline Car avec la soif, la faim, plus rien ne fait corps.
\begin{spacing}{0.8}
~
\end{spacing}
\noindent À force de salir l’essence de la vie,
\newline Un point de non-retour nous a crûment surpris.e.s.
\newline Reflet douloureux, regret qui laisse des traces,
\newline Le sens dégouline, l’adversité nous enlace.
\begin{spacing}{0.8}
~
\end{spacing}
\noindent Assumer,
\newline Admettre nos torts,
\newline Endurer l’effort,
\newline Endosser.
\begin{spacing}{0.8}
~
\end{spacing}
\noindent Solitude d’une rédemption palpable,
\newline Quand d’un souvenir, le réveil sonne coupable.
\newline Elle fait mal cette responsabilité.
\newline Un temps, privilégié.e.s, maintenant, réfugié.e.s.
\begin{spacing}{0.8}
~
\end{spacing}
\noindent Injustice
\newline \noindent D’un rêve déçu.
\newline \noindent D’un passé déchu,
\newline \noindent Un supplice.

\begin{spacing}{1.6}
~
\end{spacing}
\noindent Et si une nouvelle histoire pouvait s’écrire ?
\begin{spacing}{0.8}
~
\end{spacing}
\noindent Est-ce que nous ne le pouvons pas ?
\newline Ou est-ce que nous ne le voulons pas ?
\newline Peut-être sommes-nous libres de vouloir,
\newline Ou peut-être sommes-nous captif.ive.s de devoir.
\newline Assurément, c’est un respect que nous accordons,
\newline Et spontanément, c’est un pardon que nous échangeons.
\begin{spacing}{0.8}
~
\end{spacing}
\noindent Ainsi, des larmes se rallume alors la flamme.
\newline Brûler nos peines pour arroser notre feu,
\newline C’est là une nouvelle ère que l’on promeut,
\newline Une histoire qui redémarre et qui s’exclame.
\begin{spacing}{0.8}
~
\end{spacing}
\noindent Goûtons la sagesse, croquons la simplicité.
\newline Trinquons ensemble à la vie pour boire l’espoir.
\newline Descendons l’hubris pour gravir l’humilité.
\newline À nous d’y croire, il ne sera jamais trop tard.
\begin{spacing}{0.8}
~
\end{spacing}
\noindent Nous avons enfin saisi pourquoi,
\newline Parce qu’il y a urgence…
\newline La sobriété,
\newline Toi et moi,
\newline Nous.
\begin{spacing}{0.8}
~
\end{spacing}
\noindent Mirage ou virage ?
\newline Second souffle


\end{adjustwidth}
